\documentclass[a4paper]{article}
% Español (México) con XeLaTeX: fontspec para fuentes Unicode, polyglossia
% para la división silábica y los nombres del documento en la variante mexicana.
\usepackage{fontspec}
\usepackage{polyglossia}
\setdefaultlanguage[variant=mexican]{spanish}
% Los nombres chinos van como «pinyin (original)»: xeCJK da a los caracteres
% Han su propia fuente; sin ella XeLaTeX los omite con un aviso.
% FandolSong no cubre algunos caracteres de nombres (U+73FA); WenQuanYi Zen Hei sí.
\usepackage[AutoFallBack=true]{xeCJK}
\setCJKmainfont{FandolSong-Regular.otf}
\setCJKfallbackfamilyfont{\CJKrmdefault}{WenQuanYi Zen Hei}
% polyglossia no traduce los nombres de listings.
\AtBeginDocument{\providecommand{\lstlistingname}{}\renewcommand{\lstlistingname}{Listado}%
  \providecommand{\lstlistlistingname}{}\renewcommand{\lstlistlistingname}{Índice de listados}}
\usepackage{amsmath, amssymb}
\usepackage{graphicx}
\usepackage[margin=2.25cm]{geometry}
\usepackage[most]{tcolorbox}
\usepackage{booktabs}
\usepackage{tabularx}
\usepackage{array}
\usepackage{float}
\usepackage{xcolor}
\usepackage{hyperref}
\usepackage{fancyhdr}
\setCJKmainfont[AutoFakeBold=2.4]{AR PL UMing CN}
\setCJKsansfont[AutoFakeBold=2.4]{Droid Sans Fallback}
\setCJKmonofont[AutoFakeBold=2.4]{Droid Sans Fallback}
\hypersetup{colorlinks=true, linkcolor=blue!55!black, urlcolor=blue!60!black}
\graphicspath{{./}{figures/}}
\setlength{\parskip}{0.55em}
\setlength{\parindent}{2em}
\setlength{\emergencystretch}{3em}
\renewcommand{\arraystretch}{1.18}
\newtcolorbox{knowledgebox}[1]{enhanced, colback=blue!5!white, colframe=blue!70!black, colbacktitle=blue!70!black, coltitle=white, fonttitle=\bfseries, title={#1}, sharp corners, breakable}
\newtcolorbox{importantbox}[1]{enhanced, colback=yellow!10!white, colframe=yellow!75!black, colbacktitle=yellow!75!black, coltitle=black, fonttitle=\bfseries, title={#1}, sharp corners, breakable}
\newtcolorbox{warningbox}[1]{enhanced, colback=red!5!white, colframe=red!70!black, colbacktitle=red!70!black, coltitle=white, fonttitle=\bfseries, title={#1}, sharp corners, breakable}
\pagestyle{fancy}
\fancyhf{}
\fancyfoot[C]{\thepage}
\fancyfoot[R]{\footnotesize\color{gray}\href{https://github.com/hqhq1025/ai-course-notes}{AI Course Notes}}
\renewcommand{\headrulewidth}{0pt}
\renewcommand{\footrulewidth}{0pt}
\newcommand{\videourl}{https://www.youtube.com/watch?v=YshXmh_q_Q4}
\newcommand{\repourl}{https://github.com/hqhq1025/ai-course-notes}

\begin{document}
\begin{titlepage}
\centering
{\Large Notas didácticas del informe trimestral global sobre modelos grandes\par}
\vspace{1.0cm}
{\huge\bfseries 2025 Q1: el regreso del preentrenamiento, entornos de Coding, el Agent como nueva especie y la línea principal hacia la AGI}\par
\vspace{0.5cm}
{\Large Zhang Xiaojun Podcast: Zhang Xiaojun (张小珺) conversa con Li Guangmi (李广密)\par}
\vspace{0.35cm}
{\large Elaborado a partir del video público, la estructura de capítulos, la transcripción, la descripción del video y diagramas conceptuales propios\par}
\vspace{0.3cm}
{\large 2025-03-30\par}
\vspace{0.85cm}
\includegraphics[width=0.88\textwidth,height=0.42\textheight,keepaspectratio]{cover.jpg}\par
\vfill
\begin{tcolorbox}[width=0.92\textwidth, colback=black!2!white, colframe=black!55, sharp corners]
\textbf{Título del video}: 97. Informe trimestral de modelos grandes del Q1 de 2025: conversación con Guangmi sobre el mayor disenso actual, la línea principal y la cumbre principal de la AGI\par
\textbf{Canal del video}: Zhang Xiaojun Podcast\par
\textbf{Duración del video}: 02:01:11\par
\textbf{Enlace del video}: \href{\videourl}{\nolinkurl{\videourl}}\par
\textbf{Nota sobre las fuentes}: YouTube no ofrece subtítulos descargables; el texto se elaboró a partir de una transcripción por bloques con faster-whisper large-v3 (CPU, int8), la descripción del video, la estructura de capítulos y figuras didácticas propias. Las tomas fijas del podcast no se reutilizan en el texto.\par
\textbf{Página del proyecto}: \href{\repourl}{\nolinkurl{\repourl}}\par
\end{tcolorbox}
\end{titlepage}

\begingroup
\small
\setlength{\parskip}{0pt}
\renewcommand{\baselinestretch}{0.92}\selectfont
\tableofcontents
\endgroup
\newpage

\section{Introducción: ¿cuál es el hilo conductor de este informe trimestral?}

Esta sección establece primero el marco de lectura de todo el episodio. El «Informe trimestral global de grandes modelos» de Zhang Xiaojun (张小珺) no es un recuento de noticias una por una, sino que usa un conjunto de juicios para condensar los cambios del trimestre en una hoja de ruta tecnológica. El cambio clave del primer trimestre de 2025 es que Li Guangmi (李广密) vuelve a subrayar el papel fundamental del Pre-training y, al mismo tiempo, discute Coding, Agent, Online Learning, la divergencia estratégica de las empresas de modelos y el panorama entre China y Estados Unidos dentro de un mismo hilo conductor hacia la AGI.

\begin{importantbox}{Tesis central de este episodio}
El juicio central de Guangmi puede condensarse en una sola frase: el aumento de la inteligencia sigue siendo el único hilo conductor; el Pre-training se encarga de abrir el límite superior intrínseco del modelo, Coding ofrece el entorno de acción digital más general, Agent es una nueva especie y Online Learning podría ser la siguiente ruta de nivel paradigmático; tanto las empresas de producto como las empresas de modelos deben reordenar su posición en torno a estas capacidades.
\end{importantbox}

\begin{knowledgebox}{Nota sobre la estrategia visual}
Este video es una toma fija de pódcast, sin proyección de pantalla, pizarrón ni demostración de producto. En el texto, la portada se usa solo para identificar la fuente; todas las imágenes del texto son diagramas conceptuales propios, que sirven para explicar rutas tecnológicas, estrategias organizacionales, barreras de producto y el panorama de la industria. Esto se ajusta mejor al criterio de este repositorio para el tratamiento de pódcasts que repetir tomas del conductor o del invitado.
\end{knowledgebox}

\begin{figure}[H]
\centering
\includegraphics[width=0.92\textwidth]{q1-model-map.png}
\caption{Mapa del informe trimestral de grandes modelos del primer trimestre de 2025: Pre-training, Coding, Agent, Online Learning y el panorama global conforman juntos el hilo conductor. Diagrama conceptual propio, elaborado a partir de 00:00:08--00:04:22 y del contenido de todo el video.}
\end{figure}

\begin{knowledgebox}{Lectura de la figura: no leer el informe trimestral como una lista de noticias}
En la figura, Q1 2025 está en el centro, rodeado por Pre-training, Coding, Agent, Online, producto y panorama. Al leer esta figura, observe primero las relaciones de dependencia: la capacidad del modelo determina el límite superior utilizable, Coding y Agent ofrecen el entorno de acción, los productos recogen los beneficios de la inteligencia y el panorama de las empresas refleja cómo apuestan las organizaciones por estas rutas.
\end{knowledgebox}

\subsection{Ruta de lectura}

Esta sección presenta la ruta de lectura. Todo el video puede desglosarse en seis preguntas: primera, ¿por qué, cuando todos se orientan hacia las aplicaciones y el post-entrenamiento, Guangmi vuelve a subrayar el Pre-training? Segunda, ¿por qué Coding no es solo programación, sino el entorno más importante en la etapa temprana de la AGI? Tercera, ¿qué problema organizacional revela la divergencia estratégica entre OpenAI y Anthropic? Cuarta, ¿por qué se dice que Agent es una nueva especie? Quinta, ¿puede Online Learning cambiar el paradigma de entrenamiento? Sexta, ¿cómo se restringen mutuamente los modelos, los productos, el capital y el panorama entre China y Estados Unidos?

\noindent\begin{tabularx}{\textwidth}{p{0.19\textwidth}p{0.34\textwidth}X}
\toprule
Pregunta de lectura & Material de la entrevista & Juicio que hay que formar \\
\midrule
¿De dónde viene el límite superior del modelo? & Pre-training, Post-training, RL, datos sintéticos & El post-entrenamiento puede dar forma, pero el límite superior de la base sigue proviniendo del preentrenamiento y de los datos, el cómputo y la arquitectura. \\
¿Cómo se lleva Agent a la práctica? & Coding, computer use, tool use, long context & El entorno digital ofrece retroalimentación verificable; es donde Agent llega primero a la práctica. \\
¿Por qué divergen las empresas? & OpenAI, Anthropic, DeepSeek, Manus, Cursor & La estrategia no es un eslogan, sino la expresión de las capacidades organizacionales, los puntos de entrada de producto y la acumulación tecnológica. \\
¿Dónde está la barrera del producto? & Cloud, OS, contenedores, costo de token, «ladrones del fuego» & En la era del modelo desnudo se debilita; se fortalecen los entornos y flujos de trabajo que recogen la inteligencia. \\
¿Cuál es el siguiente paradigma? & Online Learning, long-term memory, retroalimentación del entorno & Si el modelo puede aprender mientras actúa, la ruta de aumento de la inteligencia volverá a cambiar. \\
\bottomrule
\end{tabularx}

\begin{warningbox}{Límite importante: este es un informe trimestral de opinión, no un anuario de hechos}
Esta nota organiza las opiniones del invitado en un marco de análisis, pero no presenta como hechos juicios del tipo «la AGI se logrará en dos años», «los pesos de la cartera de cierta empresa» o «cierta empresa fracasará». Otro límite terminológico: el Perplexity que aparece en este episodio se refiere al producto o empresa de búsqueda con IA, no a la perplexity de la evaluación de modelos de lenguaje (PPL, el indicador de perplejidad que se obtiene al aplicar la exponencial a la entropía cruzada y que intuitivamente puede entenderse como el número efectivo de candidatos por token).
\end{warningbox}

\subsection{Resumen de la sección}

El valor técnico de este informe trimestral radica en que reorganiza varias líneas de interés del primer trimestre de 2025 en un solo hilo conductor hacia la AGI: el límite superior del modelo, el entorno de acción digital, la conversión de Agent en producto, el paradigma de aprendizaje en línea, la estrategia de las empresas y las restricciones geopolíticas. Las secciones siguientes lo desarrollan nivel por nivel.
\section{El regreso del relato del Pre-training}

La sección anterior expuso el hilo conductor de todo el episodio; esta sección aborda la primera postura no consensuada. Después de 2024, la discusión en la industria se desplazó en gran medida hacia el Post-training, el RL, los reasoning model de la serie o, las aplicaciones de Agent y la conversión en producto, por lo que se volvió común el juicio de que «los rendimientos del Pre-training se están desacelerando». Aquí Guangmi (广密) sostiene lo contrario: el Pre-training no ha terminado y sigue siendo la etapa más esencial para elevar el límite superior de nuevas capacidades.

\begin{figure}[H]
\centering
\includegraphics[width=0.92\textwidth]{pretrain-posttrain-split.png}
\caption{Pre-training y Post-training: el preentrenamiento determina el límite superior intrínseco; el posentrenamiento y el RL se parecen más a dar forma a las capacidades. Diagrama conceptual propio, elaborado a partir de la conversación entre 00:04:22 y 00:11:56.}
\end{figure}

\begin{knowledgebox}{Lectura de la figura: el límite superior y la forma de las capacidades no son el mismo problema}
A la izquierda, el Pre-training se encarga de comprimir a gran escala el conocimiento del mundo, el código, los patrones de razonamiento y las capacidades latentes; a la derecha, el Post-training/RL se encarga de alinear esas capacidades con tareas, preferencias, herramientas y límites de seguridad. El juicio de Guangmi es que el posentrenamiento puede hacer que el modelo «use mejor las capacidades que ya tiene», pero difícilmente puede crear de la nada capacidades de las que carece el modelo base.
\end{knowledgebox}

\subsection{Por qué el Pre-training sigue siendo una postura no consensuada}

Esta sección desglosa de dónde surge lo «no consensuado». Del lado del consenso, mucha gente observó que después de GPT-4 el ritmo de lanzamiento de modelos se hizo más lento, aumentaron las discusiones sobre el muro de datos y los benchmark de RL/reasoning subieron de puntuación con rapidez, de modo que concluyó que los beneficios del preentrenamiento del modelo base están por agotarse. Del lado no consensuado, Guangmi considera que la desaceleración temporal de un líder no equivale al fin de un paradigma; si el equipo de Pre-training de OpenAI atraviesa inestabilidad, desde fuera podría confundirse la situación organizacional de una empresa con el límite técnico de toda la línea de desarrollo.

\begin{importantbox}{Guangmi enfatiza: el líder no es la línea de desarrollo misma}
Cuando una empresa líder se desacelera en cierta línea, la industria tiende a interpretarlo como que «esta línea técnica se terminó». Pero también puede tratarse solo de cambios en la capacidad organizacional, la rotación de talento, la asignación de recursos y las prioridades estratégicas. Para juzgar si una línea técnica ha terminado, hay que observar el desempeño conjunto de varias empresas, varias arquitecturas y la siguiente generación de modelos.
\end{importantbox}

\noindent\begin{tabularx}{\textwidth}{p{0.20\textwidth}p{0.34\textwidth}X}
\toprule
Nivel de análisis & Dónde es fácil equivocarse & Lectura más prudente \\
\midrule
Avance de una empresa & El Pre-training de OpenAI se desacelera en cierta etapa & Puede ser un problema organizacional y estratégico; no equivale a que toda la industria se desacelere. \\
Benchmark & Los reasoning model suben de puntuación con rapidez & Las puntuaciones suben rápido, pero eso no necesariamente eleva todo el límite superior de capacidades. \\
Muro de datos & Hay cada vez menos texto de internet de alta calidad & Los datos sintéticos, los entornos de código y la retroalimentación de herramientas pueden cambiar la oferta de datos. \\
Interés por las aplicaciones & Los Agent y las aplicaciones atraen más atención & El auge de las aplicaciones sigue dependiendo de que la inteligencia del modelo base siga empujando hacia arriba. \\
\bottomrule
\end{tabularx}

\subsection{Qué puede hacer y qué no puede hacer el Post-training y el RL}

Esta sección aclara el papel del Post-training/RL. Guangmi no niega el posentrenamiento; se opone a tratarlo como sustituto de las capacidades del modelo base. El posentrenamiento puede hacer que el modelo siga mejor las instrucciones, sea más seguro, maneje mejor herramientas específicas y razone mejor en tareas de matemáticas y código; el RL puede, mediante señales de recompensa, activar patrones de razonamiento que el modelo ya posee. Pero si el conocimiento, las representaciones y los priores de acción del base model son insuficientes, el posentrenamiento se parece a «un alumno de primaria que resuelve ejercicios en serie»: la calificación sube a corto plazo, pero el techo a largo plazo es limitado.

\begin{knowledgebox}{Términos clave: Pre-training, Post-training, RL}
\noindent\begin{tabularx}{\textwidth}{p{0.18\textwidth}p{0.36\textwidth}X}
\toprule
Término & Mecanismo & Papel en este episodio \\
\midrule
Pre-training & Entrenar el base model con datos a gran escala para aprender representaciones generales y capacidad de predicción & Determina el límite superior intrínseco del modelo; es el hilo conductor que Guangmi vuelve a enfatizar. \\
Post-training & Dar forma al comportamiento del modelo con datos de instrucciones, preferencias, seguridad y tareas & Mejora la usabilidad, pero sobre todo activa y organiza capacidades existentes. \\
RL & Optimizar el comportamiento mediante señales de recompensa; se usa con frecuencia en reasoning, código y tareas con herramientas & Puede reforzar los patrones de exploración y resolución de problemas, pero depende de las capacidades del modelo base. \\
Synthetic Data & Datos de entrenamiento generados por un modelo o un sistema & Podría aliviar el muro de datos, en especial las trayectorias de CoT/herramientas de alto valor. \\
\bottomrule
\end{tabularx}
\end{knowledgebox}

\begin{warningbox}{Error común: confundir «rendir mejor en exámenes» con «ser más inteligente»}
Los reasoning model avanzan rápido en matemáticas, código y benchmark, pero eso no significa que ya hayan adquirido todas las capacidades nuevas. Hay que distinguir entre «aprovechar mejor capacidades existentes» y «que las representaciones subyacentes realmente mejoren».
\end{warningbox}

\subsection{Las dificultades de los datos sintéticos y de los marcos de entrenamiento}

Esta sección explica con más detalle por qué el Pre-training todavía puede continuar. Guangmi menciona que los datos de CoT de alto valor, los datos generados por RL y el sampling dentro de entornos podrían volver a incorporarse a la etapa de preentrenamiento y aliviar el data bottleneck. Pero no se trata simplemente de volver a meter el texto generado en el conjunto de entrenamiento: exige que el marco de entrenamiento integre de forma más estrecha inference, sampling, reward, filtering y pre-training. Dicho de otro modo, el preentrenamiento del futuro podría dejar de ser un entrenamiento sobre un corpus estático y convertirse en una fábrica de datos con exploración, generación y retroalimentación.

\begin{importantbox}{Experiencia práctica: el muro de datos no se resuelve solo con «más páginas web»}
Si en el futuro los datos de alto valor provienen de la exploración del modelo, la ejecución de código, las llamadas a herramientas y la retroalimentación del entorno, la ingeniería de datos pasará de «limpiar texto de internet» a «ejecutar entornos, muestrear trayectorias, evaluar resultados y reincorporarlos al entrenamiento». Esa es también la razón por la que los Agent y el Online Learning regresan al hilo central del entrenamiento de modelos.
\end{importantbox}

\subsection{Resumen de la sección}

El regreso del relato del Pre-training no significa que el posentrenamiento carezca de importancia, sino que el posentrenamiento no puede sustituir el límite superior del modelo base. La pregunta decisiva del primer trimestre de 2025 es qué empresas todavía pueden seguir ampliando las capacidades del base model y cuáles solo están empaquetando capacidades existentes en productos y benchmark de mejor apariencia.
\section{Coding: el entorno más general del mundo cibernético}

La sección anterior trató el límite superior de los modelos; esta sección pasa al entorno de acción. Guangmi (广密) llama a Coding «el entorno del mundo cibernético más general» y «las manos del modelo». La frase no afirma que la AGI equivalga a escribir código, sino que el entorno de código tiene tres ventajas: se puede operar, se puede verificar y se puede componer. Cuando el modelo escribe código, lo ejecuta, lee la retroalimentación de las pruebas, llama herramientas y modifica archivos, en esencia está actuando dentro de un mundo digital controlable.

\begin{figure}[H]
\centering
\includegraphics[width=0.96\textwidth]{coding-as-cyber-environment.png}
\caption{Coding es un entorno del mundo cibernético: el código no es una habilidad aislada, sino la puerta de entrada para que el modelo actúe, verifique y transforme el mundo digital. Diagrama conceptual de elaboración propia, a partir de la conversación de 00:11:56--00:19:55.}
\end{figure}

\begin{knowledgebox}{Lectura de la figura: el valor de Coding está en el entorno, no en la sintaxis}
De Code a Tools, Tests, Agent y AGI Path, la figura destaca el ciclo cerrado: el modelo puede producir acciones, las acciones pueden cambiar el entorno y el entorno puede devolver retroalimentación. Frente a las preguntas y respuestas de un chat, en el entorno de código es más fácil definir el éxito y el fracaso; por eso resulta más adecuado para entrenar y evaluar un Agent.
\end{knowledgebox}

\subsection{Por qué Coding es más fundamental que la búsqueda y la recomendación}

Esta subsección explica la analogía de Guangmi. El motor de búsqueda organiza las páginas web como un entorno consultable; el motor de recomendación organiza el contenido como un entorno consumible; Coding, en cambio, organiza la actividad económica digital como un entorno modificable. El modelo no solo encuentra información: también puede crear archivos, llamar API, ejecutar programas, escribir scripts de automatización, probar resultados y corregir errores. Esta propiedad de «poder cambiar el mundo» hace que Coding esté más cerca del cuerpo de un Agent que la simple distribución de información.

\begin{knowledgebox}{Nota de clase: el código es el espacio de acciones del modelo en el mundo digital}
En el lenguaje del aprendizaje por refuerzo, el entorno proporciona estados, las acciones cambian los estados y la retroalimentación evalúa las acciones. El entorno de Coding vuelve muy claros estos elementos: el repo, los archivos, las pruebas, los mensajes de error, la CI, las API, el navegador y la base de datos pueden funcionar como estado y como retroalimentación.
\end{knowledgebox}

\noindent\begin{tabularx}{\textwidth}{p{0.18\textwidth}p{0.34\textwidth}X}
\toprule
Entorno & Acciones principales & Valor para un Agent \\
\midrule
Búsqueda & Consultar, filtrar, resumir & Proporciona información, pero su capacidad de cambiar el mundo es débil. \\
Recomendación & Emparejar, ordenar, consumir & Se orienta más a distribuir tráfico; no es un entorno de acción general. \\
Coding & Leer y escribir código, ejecutar pruebas, llamar herramientas, corregir errores & Se puede operar, verificar y componer; es un entorno de acción digital. \\
Science/Gaming & Experimentar, simular, buscar estrategias & También es sólido, pero por ahora su generalidad y su escala de datos son menores que las de Coding. \\
\bottomrule
\end{tabularx}

\subsection{Computer use y las manos del modelo}

Esta subsección conecta Coding con computer use. Guangmi menciona que Manus hizo que el público percibiera con fuerza, por primera vez, el magic moment de tool use/computer use, y que la capacidad de los modelos de Anthropic es una base importante detrás de ello. Aquí conviene notar una relación de niveles: una empresa de producto puede envolver una capacidad en una experiencia perceptible, pero el modelo subyacente debe tener primero la capacidad de operar herramientas, entender la UI, leer y escribir archivos y mantener el estado en tareas largas.

\begin{importantbox}{La línea divisoria entre «saber decir» y «saber hacer»}
En la era de ChatGPT, los modelos hacían sentir sobre todo que «entienden»; en la era de los Agent, los modelos deben hacer sentir que «pueden completar la tarea». Este cambio requiere herramientas, permisos, memoria, recuperación ante errores y verificación de resultados, no solo respuestas más fluidas.
\end{importantbox}

\subsection{¿Coding se convertirá en la forma final del producto?}

Esta subsección responde a la pregunta de seguimiento de Zhang Xiaojun (张小珺). El juicio de Guangmi se acerca más a esto: Coding es un motor tecnológico, no necesariamente la expresión final del producto. El producto final de un motor de recomendación puede ser un flujo de videos cortos; el producto final de un motor de búsqueda es la caja de búsqueda y la página de respuestas; Coding, como motor, podría manifestarse en el futuro en IDE, automatización de oficina, Agent personales, flujos de trabajo empresariales, cómputo científico y fábricas de software. Para el usuario común, la interfaz final quizá diluya la «programación», pero por dentro seguirá habiendo ejecución de código y de herramientas.

\begin{warningbox}{No hay que equiparar el mercado de herramientas para desarrolladores con el mercado del motor de Coding}
Cursor, Claude Code y los plugins de IDE fueron las primeras formas en despegar, porque los desarrolladores tienen una motivación alta y una retroalimentación clara. Pero si Coding es de verdad el motor general de la economía digital, su producto final no servirá solo a programadores, sino que llegará a la oficina, la investigación, la operación, el diseño, el análisis de negocio y la automatización personal.
\end{warningbox}

\subsection{Resumen de la sección}

La importancia estratégica de Coding está en que da al modelo un mundo digital en el que se puede actuar, verificar y escalar. Es el campo de batalla principal donde los Agent se implementan primero, y también el entorno decisivo para comprobar si un modelo pasa de conversar a actuar.
\section{OpenAI vs Anthropic: la estrategia como expresión de las capacidades organizacionales}

La sección anterior explicó por qué importa Coding; esta sección examina la estrategia de las empresas. Guangmi (广密) interpreta la divergencia entre OpenAI y Anthropic como «la expresión de capacidades organizacionales distintas». Ambas empresas comparten el mismo origen, pero sus apuestas en el primer trimestre (Q1) de 2025 fueron claramente distintas: OpenAI apostó con más fuerza por la serie o de reasoning models, por ChatGPT como punto de acceso para el consumidor y por sus ambiciones de plataforma; Anthropic se concentró más en el base model, Coding, Agentic workflow y el mercado empresarial.

\begin{figure}[H]
\centering
\includegraphics[width=0.92\textwidth]{openai-anthropic-strategy.png}
\caption{OpenAI vs Anthropic: la diferencia estratégica expresa capacidades organizacionales, valores y rutas de producto. Diagrama conceptual propio, elaborado a partir de la conversación entre 00:19:55 y 00:30:18.}
\end{figure}

\begin{knowledgebox}{Lectura de la figura: la estrategia de una empresa no es un PPT, sino quién es capaz de hacer qué}
A la izquierda, OpenAI pone el acento en los modelos, el punto de acceso para el consumidor y la serie o; a la derecha, Anthropic pone el acento en la seguridad, el sector empresarial, Coding y los flujos de trabajo Agentic. El análisis de Guangmi no es una simple clasificación, sino una pregunta: ¿de qué manera el talento central de la organización, la presión del producto y las prioridades de la dirección modifican la ruta tecnológica?
\end{knowledgebox}

\subsection{La serie o, el punto de acceso para el consumidor y la competencia por recursos con el preentrenamiento}

Este apartado desglosa las tensiones internas de OpenAI. La serie o avanza con rapidez en los benchmarks de matemáticas y de código, por lo que obtiene con facilidad recursos y atención dentro de la organización; el crecimiento de usuarios de ChatGPT y su orientación hacia el internet de consumo también consumen la energía de la dirección. Al mismo tiempo, la rotación de talento en el equipo de preentrenamiento y los ajustes organizacionales podrían debilitar el ritmo de entrenamiento del modelo base. Lo que preocupa a Guangmi es que, si OpenAI se convierte demasiado pronto en una empresa de internet de consumo, la línea principal del «techo de inteligencia» se desvíe por el tráfico y el crecimiento del producto.

\begin{warningbox}{Nota de clase: el tráfico es importante, pero no equivale a la línea principal de la inteligencia}
Si una empresa de modelos se deja arrastrar demasiado pronto por el punto de acceso para el consumidor, puede obtener ventajas en crecimiento de usuarios, contenido, marca y ecosistema, pero también puede reducir su concentración en el entrenamiento del modelo base, la investigación de largo plazo y el techo de capacidades. Este juicio requiere observación continua y no puede extraerse de la popularidad de un producto en un solo trimestre.
\end{warningbox}

\subsection{La oportunidad de Anthropic en Coding}

Este apartado examina a Anthropic. La confianza de Guangmi en Anthropic se basa en dos aspectos: primero, es posible que conserve una fuerte capacidad organizacional en el base model de preentrenamiento; segundo, Claude/Sonnet ya ha recibido el voto de los desarrolladores en Coding y en Agentic workflow. Productos como Cursor invocan de forma masiva los modelos de la serie Claude, lo que muestra que Coding no es una capacidad marginal, sino un escenario capaz de generar consumo real de tokens y pagos comerciales.

\begin{importantbox}{Experiencia práctica: el voto de los desarrolladores pesa más que los eslóganes}
Si los desarrolladores están dispuestos a pagar de forma sostenida por la capacidad de Coding de un modelo, y las herramientas lo establecen como modelo predeterminado, eso indica que la capacidad del modelo ya forma parte de una cadena de productividad real. El valor del escenario de Coding no está solo en los benchmarks, sino en el código real, los errores reales y las entregas reales.
\end{importantbox}

\subsection{La divergencia de visiones en Silicon Valley: ¿importa más la inteligencia o el tráfico?}

Este apartado eleva la divergencia entre empresas a una divergencia de visiones. Un bando sostiene que el techo de inteligencia sigue siendo la raíz de todo y que hay que seguir apostando por el base model, el preentrenamiento, Coding, Agent y AI for Science; el otro bando da más peso al punto de acceso para el consumidor, el crecimiento de usuarios, el punto de acceso predeterminado y el ecosistema de plataforma. En la práctica no se trata de elegir uno u otro, pero cada empresa oscila entre ambos extremos según su organización, su capital y la presión del producto.

\noindent\begin{tabularx}{\textwidth}{p{0.20\textwidth}p{0.34\textwidth}X}
\toprule
Ruta & Indicadores de interés & Riesgos \\
\midrule
Prioridad a la inteligencia & base model, Coding, Agent, capacidad científica, techo de largo plazo & El producto y los ingresos de corto plazo pueden avanzar con lentitud. \\
Prioridad al tráfico & escala de usuarios, punto de acceso, productos de consumo, opción predeterminada & Puede diluir la atención dedicada a la investigación. \\
Prioridad a la plataforma & API, protocolos de herramientas, ecosistema, flujos de trabajo empresariales & Exige mantener a la vez la capacidad y la confianza de los desarrolladores. \\
\bottomrule
\end{tabularx}

\subsection{Resumen de la sección}

La divergencia entre OpenAI y Anthropic muestra que la competencia entre empresas de IA no es una simple competencia de parámetros de modelos, sino la expresión conjunta de capacidades organizacionales, convicciones sobre la ruta a seguir, puntos de acceso de producto, estructura de clientes y presión comercial. El AGI roadmap que sigue debe leerse en este contexto.
\section{Hoja de ruta hacia la AGI: del pie de montaña de ChatGPT a la cumbre principal de la acción}

Las secciones anteriores trataron, respectivamente, el techo de los modelos, los entornos de acción y la diferenciación entre empresas. Esta sección entra en la hoja de ruta de Guangmi (广密) como «fundamentalista de la AGI». Él sostiene que el aumento de la inteligencia es la única línea principal y que la inteligencia misma es la mayor aplicación. ChatGPT es solo el pie de la montaña; detrás vienen varias cimas más: Coding, Coding Agent, General Agent, AI for Science, Robotics, entre otras.

\begin{figure}[H]
\centering
\includegraphics[width=0.96\textwidth]{agi-roadmap-mountain.png}
\caption{La cordillera de la ruta hacia la AGI: ChatGPT es solo el pie de la montaña; las cumbres principales que siguen entran, una tras otra, en entornos de acción más exigentes. Diagrama conceptual de elaboración propia, a partir de la conversación de 00:30:18--00:43:00.}
\end{figure}

\begin{knowledgebox}{Lectura de la figura: la cumbre principal pasa del diálogo a la acción}
De ChatGPT a Coding, Coding Agent, General Agent, AI4Science y Robotics, la ruta se acerca cada vez más a la acción real y a la retroalimentación real. La conversación demuestra que el modelo puede expresar inteligencia; los Agent y la ciencia demuestran que el modelo puede usar la inteligencia para modificar su entorno.
\end{knowledgebox}

\subsection{La inteligencia misma es la mayor aplicación}

Esta subsección explica la idea de que «la inteligencia misma es la mayor aplicación». La frase no niega los productos concretos; subraya que todos los productos aprovechan el excedente de capacidad de los modelos. Productos como ChatGPT, Cursor, Manus y Perplexity surgieron porque la capacidad de los modelos base se volvió de pronto suficientemente alta, y las empresas de producto la convirtieron, mediante un entorno o un contenedor, en un flujo de trabajo que el usuario puede percibir. Las empresas de aplicaciones no crean inteligencia de la nada: en el momento oportuno, recogen los beneficios del excedente de la investigación.

\begin{importantbox}{La tarea de las empresas de aplicaciones: construir contenedores de inteligencia}
Si el modelo base es el nivel del agua que sube, la tarea de las empresas de aplicaciones no es fingir que crearon el mar, sino construir bien el puerto, las tuberías y las llaves de agua. Cuanto más cerca esté el contenedor de las tareas reales, más clara sea la retroalimentación y más controlable el costo, más posibilidades tendrá de consolidar barreras de entrada.
\end{importantbox}

\subsection{El lugar de AI for Science y Robotics}

Esta subsección examina el tramo final de la hoja de ruta. Guangmi considera que en 2026/2027 podría darse una explosión de AI for Science, porque los problemas científicos tienen alto valor, retroalimentación verificable y un enorme espacio de mercado. Robotics está más cerca del mundo físico, pero él se mantiene escéptico ante el enfoque actual de los robotics foundation model: si solo se monta un VLM/VLA sobre un robot sin resolver de verdad los problemas de acción, datos, retroalimentación y cuerpo físico, el enfoque no llega a lo esencial.

\begin{warningbox}{No hay que mitificar Robotics antes de tiempo}
Los robots son un destino final importante en la ruta hacia la AGI, pero sus datos, su hardware, su retroalimentación física y sus restricciones de seguridad son mucho más complejos que los del mundo digital. Frente a obtener primero retroalimentación estable en los entornos de Coding y de Agent, apostar directamente por el mundo físico puede resultar más lento, más caro y más difícil de verificar.
\end{warningbox}

\subsection{Cortina de humo y línea principal: dónde ubicar la generación de imágenes a partir de texto}

Esta subsección aborda un juicio que se presta a controversia. Guangmi mencionó que la generación de imágenes a partir de texto podría ser una cortina de humo de OpenAI. No quiere decir que la generación de imágenes carezca de valor, sino que puede atraer demasiada atención sin representar necesariamente la máxima prioridad de la línea principal de la inteligencia. Para un fundamentalista de la AGI, la generación de imágenes, el video corto y los éxitos de consumo masivo deben volver a la misma pregunta: ¿hace que el modelo sea más inteligente, más capaz de actuar y más capaz de aprender?

\begin{knowledgebox}{Nota de clase: filtrar los temas de moda en IA según si aumentan la capacidad de acción}
Un tema de moda en IA puede ser muy comercial, muy llamativo y muy viral, pero si no profundiza la capacidad del modelo para entender el mundo, operar herramientas, obtener retroalimentación o transferir lo aprendido, no necesariamente pertenece a la línea principal de la AGI. El valor del informe trimestral está precisamente en usar la línea principal para filtrar los temas de moda.
\end{knowledgebox}

\subsection{Resumen de la sección}

El AGI roadmap de Guangmi sitúa a ChatGPT al pie de la montaña y coloca los entornos de acción y el descubrimiento científico en las cumbres principales que siguen. Esta hoja de ruta no necesariamente predice bien los plazos, pero ofrece un marco de juicio: cuanto más permita algo que el modelo actúe, verifique, aprenda y modifique el mundo, más cerca está de la línea principal.
\section{Cómo medir la inteligencia: supervivencia, exploración, automatización y economía de los token}

La sección anterior trazó una hoja de ruta; esta sección responde a la pregunta insistente de Zhang Xiaojun (张小珺): «¿cuál es la esencia de la inteligencia?». Guangmi (广密) no dio una definición rigurosa, sino que describió la evolución de la inteligencia humana con tres palabras clave: supervivencia, exploración y automatización. Aplicado a los modelos, el progreso de la inteligencia puede entenderse como la capacidad del modelo de adaptarse mejor a su entorno, descubrir nuevas estrategias, completar tareas de forma automática y realizar trabajos más complejos con un costo en token menor o más eficiente.

\begin{figure}[H]
\centering
\includegraphics[width=0.92\textwidth]{intelligence-measurement.png}
\caption{Cómo se mide el progreso de la inteligencia: de la supervivencia, la exploración y la automatización al consumo de token y al ciclo cerrado de tareas. Diagrama conceptual propio, elaborado a partir de la conversación de 00:43:00--00:48:03.}
\end{figure}

\begin{knowledgebox}{Lectura de la figura: la inteligencia no es una cifra de benchmark}
En la figura aparecen supervivencia, exploración, automatización, Token, tarea y retroalimentación. Al leerla conviene notar lo siguiente: un benchmark es un corte de medición, mientras que la inteligencia real se parece más a un ciclo cerrado de tareas. Que el modelo pueda explorar de forma independiente, elegir herramientas, consumir una cantidad razonable de token y completar tareas verificables es una forma de medición más cercana a la era de los Agent.
\end{knowledgebox}

\subsection{Por qué importa el consumo de token}

Esta sección explica la economía de los token. En la entrevista se menciona que una conversación ordinaria con un chatbot puede consumir algunos miles de token, que una tarea de búsqueda o de investigación puede consumir más, y que una tarea promedio de Manus puede llegar a cientos de miles de token. Las cifras no deben tomarse como estadísticas precisas; lo importante es la tendencia: las tareas de un Agent consumen más razonamiento, contexto, llamadas a herramientas y estados intermedios que las tareas de conversación. Si la fijación de precios de un producto sigue copiando la cuota mensual del SaaS, puede quedar desalineada respecto del costo real de cómputo.

\begin{importantbox}{Experiencia práctica: el precio de un producto de IA debe considerar el costo en token}
El SaaS tradicional tiene un costo marginal bajo, mientras que un Agent puede consumir en cada tarea una gran cantidad de token, de llamadas a herramientas y de cómputo externo. Si el producto se vende con una cuota mensual fija, pero los usuarios ejecutan muchas tareas de alto costo, la economía de los token terminará volviéndose en contra del modelo de negocio.
\end{importantbox}

\subsection{De long context a long-term memory}

Esta sección prepara el terreno para la discusión posterior sobre los Agent. Guangmi menciona que el siguiente milestone de la AGI podría ser la long-term memory, que sustituiría a la simple long context. La long context consiste en meter más material en la ventana actual; la long-term memory es la capacidad del modelo de acumular, seleccionar, actualizar y recuperar de forma continua la experiencia previa. La primera resuelve «ver más en esta ronda»; la segunda, «entender más y hacerlo mejor cuanto más tiempo se usa».

\begin{knowledgebox}{Términos clave: Long Context y Long-term Memory}
\noindent\begin{tabularx}{\textwidth}{p{0.22\textwidth}p{0.34\textwidth}X}
\toprule
Concepto & Significado & Función en un Agent \\
\midrule
Long Context & La ventana de una sola inferencia admite material más extenso & Permite que el modelo lea de una vez más archivos, páginas web e historial. \\
Long-term Memory & Guardar y actualizar experiencia entre sesiones y entre tareas & Permite que el modelo forme un estado persistente del usuario, del proyecto y del entorno. \\
Task State & Objetivo, pasos, errores y productos de la tarea actual & Determina si el Agent puede ejecutar tareas de largo alcance. \\
Experience Reuse & Transferir éxitos y fracasos pasados a tareas nuevas & Es una de las capacidades previas al Online Learning. \\
\bottomrule
\end{tabularx}
\end{knowledgebox}

\subsection{Resumen de la sección}

La medición de la inteligencia no puede basarse solo en la puntuación de preguntas aisladas. En la era de los Agent importan más el ciclo cerrado de tareas, la retroalimentación del entorno, el costo en token, la memoria de largo plazo y la reutilización de la experiencia. En conjunto, estas variables determinan si el modelo realmente pasa de «saber responder» a «saber hacer».
\section{El Agent es una nueva especie}

La sección anterior trató la medición de la inteligencia; esta sección pasa al Agent. Guangmi (广密) llama al Agent «nueva especie» porque ya no es solo un sistema que produce lenguaje, sino un sistema capaz de percibir el entorno, mantener el estado de la tarea, invocar herramientas, ejecutar acciones, leer la retroalimentación y seguir avanzando hacia el objetivo. Cuanto más cerca se esté de la AGI, más podría parecerse la evolución del Agent a un salto de fase que a una mejora lineal.

\begin{figure}[H]
\centering
\includegraphics[width=0.92\textwidth]{agent-three-capabilities.png}
\caption{Las tres capacidades decisivas del Agent: Long context reasoning, Tool use e Instruction following sostienen en conjunto a la nueva especie. Diagrama conceptual propio, elaborado a partir de la conversación de 00:48:03--00:55:49.}
\end{figure}

\begin{knowledgebox}{Lectura de la figura: las tres capacidades son indispensables}
Long Context permite al Agent mantener el estado de las tareas largas, Tool Use le permite modificar el mundo exterior e Instruction Following evita que se aparte del objetivo del usuario. Memory, Feedback y Autonomy son un nivel superior: el Agent no solo ejecuta la orden actual, sino que también acumula experiencia y explora por iniciativa propia.
\end{knowledgebox}

\subsection{Long context reasoning: el estado de las tareas largas}

Esta subsección examina primero el razonamiento con contexto largo. Una tarea de Agent rara vez es una sola pregunta y respuesta: abarca varios pasos, varios archivos, varias herramientas y varias rondas de recuperación de errores. El modelo debe recordar el objetivo, las restricciones, las rutas ya intentadas, las causas de los fallos, los productos actuales y el paso siguiente. El contexto largo es solo el requisito mínimo; la verdadera dificultad está en seleccionar, comprimir y actualizar el contexto.

\begin{warningbox}{El contexto largo no es una memoria universal}
Meter todo en la ventana trae costo, ruido y dilución de la atención. Un buen Agent necesita un estado de tarea estructurado, un sistema de archivos externo, recuperación de información, resúmenes y memoria a largo plazo, no acumular tokens sin límite.
\end{warningbox}

\subsection{Tool use: del lenguaje a la acción}

Esta subsección examina el uso de herramientas. La invocación de herramientas permite al modelo buscar, navegar por páginas web, leer y escribir archivos, ejecutar código, llamar a API y operar interfaces de usuario. Lleva al modelo de «generar sugerencias» a «ejecutar tareas». Pero el uso de herramientas también plantea problemas de seguridad, permisos, recuperación de errores y capacidad de auditoría. Cuanto más capaz es el Agent, más necesita límites explícitos y mecanismos de reversión.

\begin{importantbox}{El contrato básico de un producto de Agent}
El usuario da el objetivo; el Agent descompone la tarea e invoca herramientas; el sistema debe ofrecer control de permisos, visibilidad del avance, explicación de errores, toma de control humana y verificación de resultados. Sin esto, cuanto más capaz sea el Agent, mayor será el riesgo.
\end{importantbox}

\subsection{Instruction following: obediencia ante objetivos complejos}

Esta subsección explica la tercera capacidad. El Agent no solo debe saber usar herramientas, sino también ceñirse de forma sostenida al objetivo del usuario durante tareas de largo alcance: qué no puede hacer, qué recursos no puede usar, qué resultado cuenta como terminado, cuándo debe pedir confirmación y cómo resolver instrucciones contradictorias. El instruction following complejo es lo que permite al Agent pasar de una demo de juguete a un sistema en producción.

\noindent\begin{tabularx}{\textwidth}{p{0.20\textwidth}p{0.34\textwidth}X}
\toprule
Capacidad & Modo de fallo & Requisito de producto \\
\midrule
Razonamiento con contexto largo & Olvida el objetivo, repite intentos, usa información obsoleta & Estado de la tarea y memoria externa. \\
Uso de herramientas & Elige la herramienta equivocada, daña archivos, excede sus permisos & Entorno aislado (sandbox), confirmación, auditoría y reversión. \\
Seguimiento de instrucciones & Se desvía por información intermedia, ignora restricciones & Prioridades explícitas y manejo de conflictos. \\
Aprendizaje a partir de la retroalimentación & Repite el mismo error, no transfiere la experiencia & Registro de fallos, revisión retrospectiva y ciclo cerrado de evaluación. \\
\bottomrule
\end{tabularx}

\subsection{Resumen de la sección}

El Agent es una nueva especie no porque tenga un nombre nuevo, sino porque conecta el modelo de lenguaje con el entorno, las herramientas, la memoria y un ciclo cerrado de acción. La verdadera competencia pasará de la «calidad de la respuesta» a la «calidad con que se completa la tarea».
\section{Online Learning: una posible ruta de nivel paradigmático}

Las secciones anteriores explicaron por qué un Agent necesita un entorno y retroalimentación; esta sección aborda el Online Learning. Guangmi (广密) considera que, si en el futuro surge otra ruta nueva de nivel paradigmático, el Online Learning es una de las candidatas. Su núcleo no es entrenar el modelo y después desplegarlo, sino que el modelo explore de forma autónoma en un entorno, obtenga retroalimentación, actualice su experiencia e incluso, más adelante, actualice sus pesos o su memoria de largo plazo.

\begin{figure}[H]
\centering
\includegraphics[width=0.96\textwidth]{online-learning-loop.png}
\caption{Ciclo cerrado del Online Learning: el modelo explora de forma autónoma en el entorno, obtiene retroalimentación y actualiza sus capacidades. Diagrama conceptual de elaboración propia, basado en la conversación de 00:55:49--01:02:45.}
\end{figure}

\begin{knowledgebox}{Lectura de la figura: el Online Learning incorpora la producción de datos al proceso de uso}
De Environment a Action, Reward, Update y Memory/Weights, la figura subraya un ciclo continuo. El entrenamiento tradicional se parece más al aprendizaje fuera de línea; el Online Learning, en cambio, hace que el modelo aprenda mientras actúa en un entorno real o simulado. Cuantos más Agents haya, más datos de retroalimentación potenciales habrá.
\end{knowledgebox}

\subsection{Por qué el Online Learning podría ser el próximo paradigma}

Esta sección explica su atractivo. El Pre-training depende principalmente de corpus ya existentes, y el Post-training/RL depende de tareas y recompensas ya organizadas; el Online Learning, en cambio, podría generar datos nuevos a partir de las propias acciones del modelo. En entornos como Coding, videojuegos, experimentos científicos u operación de páginas web, el modelo puede intentar, fallar, corregir y sintetizar, y luego reincorporar esa experiencia. Esto desplaza el cuello de botella de los datos de «encontrar más datos antiguos» a «crear experiencia nueva de mayor valor».

\begin{importantbox}{Cambio central: del conjunto de datos al entorno}
Si el Online Learning se confirma, el centro del entrenamiento pasará de los conjuntos de datos estáticos a los entornos interactivos. Quien posea entornos, tareas, recompensas, evaluaciones y retroalimentación de usuarios reales podría tener un nuevo volante de datos.
\end{importantbox}

\subsection{Impacto en la GPU y en la narrativa de NVIDIA}

Esta sección responde a la pregunta sobre la industria que se planteó en la entrevista. El Online Learning no necesariamente debilita la demanda de GPU; por el contrario, podría cambiar su estructura: no solo se necesita entrenar, sino también muestrear, ejecutar, hacer inferencia de múltiples turnos, evaluar y reproducir a gran escala. La inferencia y el entrenamiento podrían entrelazarse más estrechamente, y la infra tendría que soportar al mismo tiempo rollout en línea, caché, simulación de entornos y auditoría de seguridad. La narrativa de NVIDIA tampoco se limita a vender tarjetas para entrenamiento, sino que se amplía en torno a la inferencia, la simulación, la robótica, los centros de datos y la pila de software.

\begin{knowledgebox}{Términos clave: infraestructura relacionada con el Online Learning}
\noindent\begin{tabularx}{\textwidth}{p{0.20\textwidth}p{0.34\textwidth}X}
\toprule
Término & Significado & Por qué importa \\
\midrule
Rollout & El Agent ejecuta una trayectoria en el entorno & Produce experiencia nueva que puede evaluarse. \\
Reward & Señal que evalúa el éxito y la calidad de la tarea & Determina qué aprenderá el modelo. \\
Replay & Reproducir trayectorias exitosas o fallidas para entrenamiento o análisis & Ayuda a aprender de los errores. \\
Simulation & Entorno controlable o mundo simulado & Reduce el costo de ensayo y error en el mundo real. \\
Memory Update & Escribir la experiencia en la memoria de largo plazo o en un estado externo & Hace que el modelo entienda mejor cuanto más se usa. \\
\bottomrule
\end{tabularx}
\end{knowledgebox}

\subsection{Riesgo: cómo evitar que el aprendizaje en línea caiga en reward hack}

Esta sección añade restricciones. El riesgo del Online Learning es que el modelo aprenda a explotar la recompensa, contamine la memoria, amplifique sesgos, eluda los límites de seguridad o aprenda estrategias erróneas a partir de retroalimentación incorrecta. Por ello, un sistema de aprendizaje en línea realmente utilizable debe contar con permisos por niveles, entornos aislados (sandbox), estados que puedan revertirse, evaluación fuera de línea, auditoría humana y filtros de seguridad.

\begin{warningbox}{El Online Learning no equivale a «dejar que el modelo aprenda lo que sea por su cuenta»}
Si el aprendizaje continuo no tiene límites, provoca contaminación de la memoria, deriva de objetivos y problemas de seguridad. Un Online Learning de alta calidad se parece más a un sistema experimental controlado: tareas claras, recompensas auditables, fallas rastreables y actualizaciones reversibles.
\end{warningbox}

\subsection{Resumen de la sección}

El potencial del Online Learning reside en que convierte el proceso de uso del Agent en un proceso de producción de datos. También trae nuevos problemas de infraestructura y de seguridad. En la siguiente etapa, quien logre operar de forma estable el ciclo cerrado de entornos, recompensas y retroalimentación podría obtener una nueva vía de mejora de los modelos.
\section{Modelo y producto: del modelo desnudo a Cloud/OS/contenedores}

La sección anterior trató el Online Learning; esta sección pasa al negocio y al producto. Guangmi (广密) considera que la era de los lanzamientos de modelos desnudos podría ir perdiendo peso y que las barreras competitivas vendrán cada vez más de Cloud, el OS, el ecosistema y los contenedores de flujo de trabajo. La razón es que otros pueden alcanzar las capacidades de un modelo, el código abierto las presiona y la API las convierte en mercancía; en cambio, el entorno donde el usuario completa sus tareas cada día y el sistema que posee el estado, los datos, los permisos y la relación de pago son lo que forma una barrera duradera.

\begin{figure}[H]
\centering
\includegraphics[width=0.92\textwidth]{model-product-moat.png}
\caption{Barreras de modelo y producto: la era de los lanzamientos de modelos desnudos pierde peso y Cloud/OS/el ecosistema se vuelven la línea de defensa. Diagrama conceptual propio, elaborado a partir de la conversación de 01:02:45--01:15:11.}
\end{figure}

\begin{knowledgebox}{Lectura de la figura: un modelo fuerte no equivale a un producto fuerte}
A la izquierda, Model destaca el techo de capacidad y la eficiencia de costos; a la derecha, Product/OS destaca el punto de entrada, el flujo de trabajo, el ecosistema, la facturación y el retorno de datos. Al leer esta figura hay que retener una idea: el modelo es el motor; el producto es la carretera, el tablero, el sistema de cobro y la experiencia de manejo.
\end{knowledgebox}

\subsection{El precio de 20 dólares y el problema de copiar el SaaS}

Esta subsección examina la fijación de precios. La entrevista pregunta por qué muchos productos de IA cuestan cerca de 20 dólares y si eso solo copia al SaaS. La cuestión implícita de Guangmi es que el costo marginal de un producto de IA es distinto del de un SaaS tradicional. Una tarea compleja de Agent puede consumir muchos token y llamadas a herramientas; si se vende como suscripción mensual barata de uso ilimitado, aparece presión sobre el margen bruto. Por eso, en el futuro, los precios podrían acercarse a un modelo mixto que combine uso, valor de la tarea, licencias empresariales por usuario y consumo de cómputo.

\begin{warningbox}{El precio no es un problema de UI, sino de economía del cómputo}
Si un producto empaqueta tareas de Agent de alto costo como uso ilimitado a bajo precio, puede crecer a corto plazo, pero a largo plazo el costo de inferencia y el abuso de los usuarios lo aplastarán. Un gerente de producto de IA debe entender los token, la caché, las llamadas a herramientas y el valor de la tarea.
\end{warningbox}

\subsection{¿Los modelos se comerán a los productos?}

Esta subsección aborda la preocupación principal de las empresas de producto. Guangmi resume la cuestión de fondo como feature system vs learning system. Las empresas de modelos pueden copiar las feature de un producto tradicional; pero si la empresa de producto acumula entorno, datos, flujos de trabajo, estado del usuario, sistema de permisos y ciclo de retroalimentación, una nueva versión del modelo no se la comerá directamente. A la inversa, si el producto solo envuelve al modelo, sin un entorno ni datos propios, el modelo base sí lo absorberá con facilidad.

\begin{importantbox}{Criterio para juzgar la barrera de una aplicación}
Para saber si una aplicación de IA tiene una barrera, no basta con mirar la interfaz y el prompt; hay que ver si posee: un punto de entrada a tareas reales, estado del usuario a largo plazo, retroalimentación verificable, flujos de trabajo propios, control de costos, retorno de datos y costo de sustitución.
\end{importantbox}

\begin{figure}[H]
\centering
\includegraphics[width=0.92\textwidth]{fire-thieves.png}
\caption{Los ladrones del fuego de los modelos: Perplexity, Cursor y Manus llevan las capacidades de los modelos de frontera a flujos de trabajo concretos. Diagrama conceptual propio, elaborado a partir de la conversación de 01:02:45--01:15:11.}
\end{figure}

\begin{knowledgebox}{Lectura de la figura: el valor de los ladrones del fuego está en el contenedor}
Perplexity lleva la capacidad del modelo a la búsqueda y la investigación; Cursor, al entorno de programación; Manus, a la ejecución general. No entrenan ellos mismos el modelo base más fuerte, sino que colocan la capacidad del modelo en un contenedor por el que el usuario está dispuesto a pagar, que usa una y otra vez y que genera retroalimentación.
\end{knowledgebox}

\subsection{Qué debe observar un inversionista}

Esta subsección convierte la cuestión del producto en un juicio de inversión. Guangmi menciona su portafolio ideal y el peso de cada empresa, pero esta nota no lo toma como recomendación de inversión. Lo más importante es el marco de análisis: en las empresas de modelos base se observan el techo de capacidad, la curva de costos y la estabilidad de la organización; en las empresas de aplicaciones, si aprovechan la ventana en que las capacidades del modelo se desbordan hacia otros usos, si acumulan flujos de trabajo y si tienen ingresos reales y costos controlables; en las empresas de infraestructura, si controlan los eslabones clave de entorno, inferencia, entrenamiento, datos y despliegue.

\begin{knowledgebox}{Nota de clase: invertir en tecnología no es apostar por palabras de moda}
Al final, la entrevista enfatiza «crear» en lugar de «moverse en el círculo». En la inversión en IA, esto significa entender las rutas tecnológicas y la cadena de valor: quién eleva la inteligencia, quién la recibe, quién proporciona el entorno, quién controla la retroalimentación y quién solo hace arbitraje de narrativa.
\end{knowledgebox}

\subsection{Resumen de la sección}

La relación entre modelo y producto no se reduce a quién se come a quién. Los modelos fuertes reducen el espacio de los productos delgados, pero también crean oportunidades para nuevos contenedores. La barrera de un producto proviene del entorno, el estado, la retroalimentación, el costo y las tareas del usuario, no de poner una simple ventana de chat encima del modelo.
\section{Panorama global de las empresas de modelos y de productos de IA}

Tras el análisis de las barreras de entrada de los productos, esta sección aborda el panorama de las empresas. En el mapa de empresas del Q1 de 2025, OpenAI, Anthropic, Google, DeepSeek, ByteDance, Mira/Thinking Machines, SSI, Cursor, Manus y otros actores aparecen en un mismo mapa competitivo. Las opiniones de Guangmi (广密) son marcadamente subjetivas, pero ayudan a entender la posición de los distintos tipos de empresa en la cadena de valor de la IA.

\begin{figure}[H]
\centering
\includegraphics[width=0.92\textwidth]{model-company-portfolio.png}
\caption{Panorama global de las empresas de modelos: cada empresa apuesta de forma distinta por el modelo, el producto, el ecosistema, el código abierto y la estructura de capital. Diagrama conceptual de elaboración propia, basado en la conversación de 01:15:11--01:54:32.}
\end{figure}

\begin{knowledgebox}{Lectura de la figura: no todas las empresas de IA compiten en la misma categoría}
Anthropic se orienta más a empresas y a Coding; OpenAI apuesta a la vez por el modelo y por el producto; ByteDance tiene productos y tráfico; DeepSeek representa la eficiencia del código abierto; Cursor y Manus forman la capa de ejecución de aplicaciones. Las condiciones de éxito de cada posición son completamente distintas, por lo que no basta con una sola clasificación de benchmark.
\end{knowledgebox}

\subsection{GPT-4.5, GPT-5 y los riesgos de OpenAI}

Esta sección repasa la conversación sobre OpenAI. La entrevista aborda si GPT-4.5 va a la delantera, por qué se retrasó GPT-5, si OpenAI corre riesgo de fracasar, su relación con Microsoft y su apoyo al protocolo MCP de Anthropic. Al organizar este material como curso, conviene verlo como un conjunto de variables: el ritmo del modelo base, el ritmo del reasoning model, el crecimiento de usuarios del producto, la estructura de la alianza en la nube, el ecosistema de protocolos y la estabilidad organizacional.

\begin{warningbox}{No juzgues a una empresa por un solo lanzamiento}
Tanto los riesgos como las ventajas de OpenAI son grandes: tiene la marca más fuerte, el principal punto de acceso de usuarios y el mejor posicionamiento de producto en la mente del público, pero también enfrenta una creciente complejidad organizacional, su relación con los proveedores de nube, el ritmo de la investigación y la presión de comercialización. Una sola versión de un modelo no decide quién gana a largo plazo.
\end{warningbox}

\subsection{DeepSeek, la eficiencia del código abierto y la señal de China}

Esta sección se ocupa de DeepSeek. La entrevista coloca a DeepSeek como la estrella del Q1 y subraya su impacto en el panorama global de los modelos y en el discurso sobre la eficiencia del código abierto. Su valor no se reduce a que «el modelo es barato»: demuestra que los equipos chinos pueden influir en la percepción global mediante la eficiencia de ingeniería, las decisiones de arquitectura, la recipe de entrenamiento y la difusión en código abierto. Guangmi incluso usa el lenguaje de un portafolio de inversión para expresar que le daría un peso elevado a DeepSeek, lo que refleja la importancia que concede a la vía de la eficiencia del código abierto.

\begin{importantbox}{El significado de la señal de DeepSeek}
DeepSeek llevó a la industria a reevaluar la posición de los equipos chinos en modelos frontier, eficiencia de entrenamiento, ecosistema de código abierto e influencia sobre los desarrolladores de todo el mundo. No es solo una historia de precios, sino una historia en la que se suman capacidad, costo, apertura y difusión.
\end{importantbox}

\subsection{Manus, Perplexity y Cursor: empresas de producto con gran capacidad de ejecución}

Esta sección se ocupa de las empresas de aplicaciones. En la entrevista, a Manus, Perplexity y Cursor se les llama «los que le roban el fuego a los modelos» e incluso, en tono de broma, «los reyes del envoltorio». La expresión parece mordaz, pero lo que en realidad destaca es su capacidad de ejecución: convierten con rapidez en experiencia de usuario las capacidades que las empresas de modelos aún no han convertido en producto. Perplexity reorganizó la búsqueda y la investigación, Cursor convirtió el Coding en un flujo de trabajo de uso frecuente y Manus llevó ante el público la experiencia de un Agent de propósito general.

\begin{knowledgebox}{Términos clave: empresas de modelos, empresas de aplicaciones y empresas de infraestructura}
\noindent\begin{tabularx}{\textwidth}{p{0.20\textwidth}p{0.34\textwidth}X}
\toprule
Tipo de empresa & Activos centrales & Riesgos principales \\
\midrule
Empresa de modelos & base model, infra de entrenamiento, talento de investigación, API/plataforma & Alto costo de entrenamiento, organización compleja, fuerte presión por el producto. \\
Empresa de aplicaciones & Punto de acceso de usuarios, flujos de trabajo, estado de las tareas, datos de retroalimentación & Que las empresas de modelos la copien o que los costos la ahoguen. \\
Empresa de infraestructura & Inferencia, entrenamiento, entornos, protocolos de herramientas, recursos en la nube & Debe adaptarse a los cambios en la demanda de los modelos y las aplicaciones. \\
Empresa de modelos de código abierto & Eficiencia de costos, comunidad, transparencia, difusión del ecosistema & Presión por la comercialización y por financiar el entrenamiento continuo. \\
\bottomrule
\end{tabularx}
\end{knowledgebox}

\subsection{MCP, protocolos y nichos del ecosistema}

Esta sección completa la capa de protocolos. La entrevista menciona que OpenAI apoya el protocolo MCP de Anthropic, lo que muestra que en el ecosistema de Agents no solo se compite por el modelo, sino también por la forma de conectar herramientas, la forma de transmitir el contexto y los límites de seguridad. Si un protocolo se convierte en estándar de facto, conectará los modelos, las herramientas, los sistemas empresariales y el ecosistema de desarrolladores. Las empresas de modelos deben evitar que el protocolo convierta su modelo en una mercancía intercambiable y, al mismo tiempo, evitar quedar aisladas del ecosistema.

\begin{importantbox}{La siguiente capa de competencia en el ecosistema de Agents son los protocolos}
Navegadores, IDE, SaaS empresariales, sistemas de archivos y bases de datos pueden convertirse en herramientas de un Agent. Quien defina cómo el modelo lee el contexto, invoca herramientas y devuelve resultados de forma segura estará definiendo las reglas de tránsito del ecosistema de Agents.
\end{importantbox}

\subsection{Resumen de la sección}

El panorama global de las empresas de IA no puede ordenarse solo según «quién tiene el modelo más fuerte». El modelo base, la eficiencia del código abierto, la ejecución de aplicaciones, el ecosistema de protocolos, las alianzas en la nube y la estructura de capital están reacomodando las posiciones. El cambio real del Q1 de 2025 es que la capacidad de los modelos empieza a dirigirse con más claridad hacia el Coding, los Agents y los productos de flujo de trabajo.
\section{El panorama China–Estados Unidos: cómo superar el bloqueo geopolítico}

La sección anterior trató el mapa de empresas; esta cierra con el panorama entre China y Estados Unidos. Al final de la entrevista se insiste en que la inversión y el emprendimiento tecnológicos no pueden basarse en moverse dentro de los círculos de contactos, sino en crear. El bloqueo geopolítico limita el cómputo, el capital, los mercados y la cooperación, pero no determina por sí solo quién gana; lo que realmente atraviesa esas restricciones es la creatividad técnica, la ejecución de ingeniería, la comprensión del producto, la difusión mediante código abierto y la capacidad de organizar la cadena industrial.

\begin{figure}[H]
\centering
\includegraphics[width=0.92\textwidth]{china-us-geofence.png}
\caption{El panorama de la IA entre China y Estados Unidos: bajo el bloqueo geopolítico, la creatividad técnica importa más que moverse en los círculos de contactos. Diagrama conceptual propio, elaborado a partir de la conversación de 01:54:32--02:01:11.}
\end{figure}

\begin{knowledgebox}{Lectura de la figura: tener ventajas distintas no obliga a limitarse a la defensa}
Del lado estadounidense están el cómputo, los frontier lab, el capital y los puntos de acceso a productos globales; del lado chino, la velocidad de ingeniería, los escenarios de aplicación, la eficiencia del código abierto y la capacidad de ejecución de la cadena industrial. El bloqueo geopolítico genera fricción, pero también obliga a los equipos a buscar nuevos caminos en eficiencia, ingeniería y ecosistema local.
\end{knowledgebox}

\subsection{Dónde están las oportunidades de los equipos chinos}

Esta sección divide las oportunidades en cuatro tipos. La primera es de eficiencia: construir modelos y productos competitivos con menos cómputo. La segunda es de aplicación: China tiene una gran cantidad de escenarios de negocio digitalizado, contenido, comercio electrónico, industria y hardware. La tercera es de código abierto: obtener distribución entre desarrolladores de todo el mundo mediante modelos y herramientas abiertos. La cuarta es de cadena industrial: cuando la IA se combina con hardware, robots, automóviles y electrónica de consumo, la cadena de suministro y la velocidad de ingeniería de China cobran importancia.

\begin{importantbox}{Experiencia práctica: en tiempos de bloqueo, hace más falta una creación verificable}
Cuanto más fuertes son las restricciones externas, menos basta con el discurso. Los equipos necesitan presentar modelos, código, productos, ingresos, curvas de eficiencia y adopción por parte de desarrolladores que puedan verificarse. La verdadera influencia global proviene de lo que se crea, no del respaldo de un círculo.
\end{importantbox}

\subsection{Riesgos bajo las restricciones geopolíticas}

Esta sección también debe mantener la conciencia de los riesgos. Los equipos chinos enfrentan limitaciones en el suministro de cómputo, el cumplimiento normativo en el extranjero, la confianza en la marca, las ventas a empresas, los pagos y el acceso a los ecosistemas; los equipos estadounidenses también enfrentan costos, regulación, burbujas de producto, expectativas del capital y una creciente complejidad organizacional. El panorama China–Estados Unidos no es una presión unilateral ni un rebase unilateral, sino dos sistemas de innovación que buscan avances bajo restricciones distintas.

\noindent\begin{tabularx}{\textwidth}{p{0.18\textwidth}p{0.34\textwidth}X}
\toprule
Dimensión & Ventajas/riesgos de EE.\,UU. & Ventajas/riesgos de China \\
\midrule
Cómputo & Suministro y ecosistema de nube sólidos, pero con alta presión de costos y regulación & Limitado por los controles de exportación; requiere eficiencia y alternativas. \\
Modelos & Alta concentración de frontier lab & La eficiencia del código abierto y las recipe de ingeniería pueden ser la vía de avance. \\
Productos & Puntos de acceso globales y clientes empresariales sólidos & Escenarios locales abundantes, pero una marca global es más difícil. \\
Cadena industrial & Ecosistema de software sólido & Cadenas sólidas de hardware, automóviles, robots y manufactura. \\
\bottomrule
\end{tabularx}

\subsection{Resumen de la sección}

Lo esencial del panorama de la IA entre China y Estados Unidos no es un simple optimismo o pesimismo, sino ver quién logra crear mejores modelos, productos, entornos y ecosistemas bajo restricciones. Para seguir a largo plazo el ámbito de la IA e internet, conviene fijarse en productos verificables y no solo en el entusiasmo que genera el discurso.
\section{Síntesis y ampliación}

Esta sección condensa todo el video en un marco reutilizable. El informe trimestral sobre modelos grandes del primer trimestre de 2025 puede entenderse como «el regreso de la inteligencia como eje central»: no hay que dejarse desviar por el éxito repentino de un solo producto ni por el lanzamiento de una sola empresa, sino seguir preguntando de dónde proviene el límite superior del modelo, cómo actúa el modelo, cómo la acción genera retroalimentación, cómo la retroalimentación vuelve al aprendizaje, cómo el producto da cabida a la inteligencia y cómo la empresa convierte su capacidad organizacional en estrategia.

\begin{importantbox}{Seis conclusiones centrales}
Primero, el preentrenamiento sigue siendo determinante para el límite superior de la capacidad base y no debe descuidarse por el auge del postentrenamiento. Segundo, Coding es el entorno digital en el que los modelos obtuvieron primero una capacidad de acción verificable. Tercero, Agent es el salto sistémico de la salida de lenguaje a la ejecución de tareas. Cuarto, Online Learning podría convertir el proceso de uso en un nuevo proceso de producción de datos. Quinto, las barreras de entrada de las aplicaciones provienen del entorno, el estado, la retroalimentación y el control de costos. Sexto, la competencia entre China y Estados Unidos se decidirá, en última instancia, por la creación verificable, y no por los relatos que circulan dentro de cada círculo.
\end{importantbox}

\subsection{Tabla de consulta rápida de términos de este episodio}

Esta sección ofrece una consulta rápida para relacionar este episodio con los episodios anteriores y posteriores de la serie de IA e internet de Zhang Xiaojun (张小珺). Hay una continuidad clara entre los productos Agent de EP97 y EP101, la ruta multimodal de EP102, la inteligencia corporizada de EP106/109, el informe técnico sobre Agent de EP110 y los informes trimestrales sobre modelos grandes de EP127/136.

\noindent\begin{tabularx}{\textwidth}{p{0.23\textwidth}p{0.35\textwidth}X}
\toprule
Término/empresa & Significado en este episodio & Cómo darle seguimiento \\
\midrule
Preentrenamiento & Abrir el límite superior intrínseco del base model & Observar si la siguiente generación de modelos muestra capacidades nuevas, y no solo puntuaciones infladas. \\
Reasoning/RL & Reforzar el razonamiento y el desempeño en tareas & Observar si se traduce en la realización de tareas reales. \\
Coding & Entorno de acción digital y «las manos» del modelo & Observar el pago de los desarrolladores, las llamadas a herramientas y la tasa de éxito en tareas de extremo a extremo. \\
Agent & Ciclo cerrado de percepción, herramientas, memoria y ejecución & Observar las tareas largas, la recuperación ante errores y la controlabilidad. \\
Online Learning & Generar retroalimentación durante el uso y aprender de forma continua & Observar si el entorno, la recompensa, la memoria y la auditoría de seguridad han madurado. \\
DeepSeek & Eficiencia de código abierto y señal de los modelos chinos & Observar la capacidad, el costo, el ecosistema y la iteración continua. \\
Manus/Cursor & Contenedores de producto para la capacidad del modelo & Observar las barreras del flujo de trabajo, los ingresos y la economía de los token. \\
\bottomrule
\end{tabularx}

\subsection{Preguntas para el seguimiento}

Esta sección convierte el informe trimestral en una lista de seguimiento. Con estas preguntas, el lector puede juzgar si el conjunto de juicios de Guangmi (广密) se sigue sosteniendo en los trimestres posteriores de 2025.

\begin{enumerate}
\item ¿La siguiente generación de base model mostrará capacidades nuevas evidentes, o dependerá sobre todo del post-training para inflar los resultados en benchmark?
\item En el escenario de Coding, ¿los modelos líderes seguirán siendo los de la serie Anthropic/Claude, o OpenAI, Google, DeepSeek y otros volverán a alcanzarlos?
\item ¿Podrán los productos Agent pasar de la demostración a la entrega estable, en especial en lo que respecta a las tareas largas, los permisos, la seguridad y el costo?
\item ¿Surgirá un paradigma de ingeniería reproducible para Online Learning, en lugar de quedarse en el plano conceptual?
\item Entre las empresas de modelos y las empresas de aplicaciones, ¿quién podrá controlar mejor el estado del usuario, el entorno de la tarea y los datos de retroalimentación?
\item ¿La ruta de eficiencia de código abierto al estilo de DeepSeek podrá seguir extendiéndose a más modelos y ecosistemas de herramientas?
\item En la competencia de IA entre China y Estados Unidos, ¿las restricciones de capacidad de cómputo se compensarán en parte con la eficiencia algorítmica, la optimización de ingeniería y el ecosistema de código abierto?
\end{enumerate}

\subsection{Lecturas adicionales}

\begin{itemize}
\item Si te interesan los productos Agent y el emprendimiento de aplicaciones, puedes contrastar con EP101 YouWare, el informe técnico sobre Agent de EP110 y la historia técnica de Agent de EP139.
\item Si te interesan la multimodalidad y los modelos del mundo, puedes contrastar con la entrevista a Zhang Xiangyu (张祥雨) de EP102 y la entrevista a Xie Saining (谢赛宁) de EP133.
\item Si te interesa la línea de la inteligencia corporizada y la robótica, puedes contrastar con EP106 Wang He (王鹤), EP109 Guanglun Zhineng (光轮智能) y la versión con proyección de pantalla sobre VLA `eiQFomOuCJs`.
\item Para la continuidad de los informes trimestrales sobre modelos grandes, puedes contrastar con EP127, EP136 y las revisiones trimestrales de 2023/2024 que se enumeran en la descripción del video.
\end{itemize}

\end{document}
